\documentclass{article}

\usepackage[margin=1in]{geometry}
\usepackage{amsmath}
\usepackage{amssymb}
\usepackage{microtype}

\title{Rewritten Chain of Thought for the Solution to the\\
Two-Minus Water-Wave Amplitude}
\author{}
\date{}

\newcommand{\pp}[1]{\left[\,#1\,\right]_{+}}

\begin{document}

\maketitle


\section*{Rewritten Chain of Thought}

I am allowed to read only \texttt{prompt.md} and \texttt{OnShellBG.m}, and I work
only from those plus data I generate myself; the deliverable goes in
\texttt{fugu\_ultra}. The prompt is unusually generous: it tells me up front that
the answer is a \emph{piecewise homogeneous polynomial} in the
$\omega_i$ --- a genuine polynomial on each chamber, with the two-minus kinematic
space cut into chambers, and a different homogeneous polynomial on each. That is
a strong hint and I take it at face value from the start: my job is not ``is it
piecewise?'' but ``what are the chamber walls and what polynomial sits on each
piece.'' Reading \texttt{OnShellBG.m}, I see immediately where the pieces must
come from: every kernel and every propagator is built out of
$\mathrm{mag}[k]=|k|$, and on this sector $k_i=\sigma_i\omega_i^2/g$, so the only
non-analytic objects in the whole construction are absolute values of signed sums
$\sum_i\sigma_i\omega_i^2$. An $|\cdot|$ has a corner where its argument changes
sign; that is exactly a chamber wall. So the hint and the code agree before I run
anything: the walls will be subset-sum surfaces.

First, is \texttt{wolframscript} available? It is. I run the file as-is, and that
is mildly wasteful: \texttt{OnShellBG.m} ends in a \texttt{VI.\ TESTS} block that
demonstrates a \emph{different} sector, $\sigma=(-1,+1,\dots,+1)$, where the
amplitude is identically zero, so it prints ``$A_5=0$, $A_6=0$, $A_7=0$'' and
then grinds on $n=8$. I let it run, then stop leaning on it. Better: split the
file as text at the \texttt{VI.\ TESTS} marker with \texttt{awk} and keep only the
definitions in a \texttt{BGdefs\_generated.m}, so I load the recursion without the
demo. (Small frictions follow all the way through: a
\texttt{Do[\dots,\{\{n,fw\},\dots\}]} trips ``Tag \texttt{List} is Protected,''
so I switch those loops to \texttt{Scan}; later a \texttt{Count[\dots,\,s\_/;\dots]}
with a bound pattern throws a syntax error and I rewrite it as
\texttt{Length@Select}.)

Now the oracle behaves, and the first real surprise is $n=4$. Direct rational
evaluation returns \texttt{Indeterminate} --- the messages are ``$1/0$''
(\texttt{Power::infy}) and ``$0\cdot\text{ComplexInfinity}$''
(\texttt{Infinity::indet}). The reason is structural and I pin it down later: at
$n=4$ the solver always returns $\omega=(-b,a,b,-a)$, so legs $2$ and $4$ carry
exactly opposite momenta, $k_2+k_4=-a^2+a^2=0$, and the BG current on $\{2,4\}$
hits a propagator $-i/(\omega_S^2/|k_S|-g)$ with $|k_S|=0$ --- a $1/0$ multiplied
by a vanishing numerator, a removable $0/0$. So I set $n=4$ aside, to recover as a
limit, and look at $n=5,6,7$ with generic positive free frequencies. Those are
finite and purely imaginary rationals: e.g.\ for free $\{2,3,4\}$ I get
$A_5=-\tfrac{8704}{3}i$, for $\{4,3,2\}$ I get $A_5=-19968\,i$, for $\{1,2,10\}$
I get $-\tfrac{2176}{13}i$. Writing $A_n=i\,P_n$ and chasing $P_n$ is the obvious
move.

For $n=4$ I do not have to guess the chamber polynomial --- I can ask Wolfram
directly. I feed symbolic free frequencies $\{a,b\}$, get the symbolic
\texttt{BGAmplitude} (a genuinely ugly thing, full of nested
$\mathrm{Abs}[a^2-b^2]$, $\mathrm{Abs}[a^2+b^2]$, and the like), and
\texttt{FullSimplify} it under sign assumptions on the two chambers:
\[
b>a:\quad A_4=-8\,i\,a^3 b,\qquad a>b:\quad A_4=-8\,i\,a\,b^3 .
\]
Running all eight sign-cells of $(a,b)$ collapses everything to just these two
polynomials, $-8ia^3b$ and $-8iab^3$, i.e.\
$A_4=-8\,i\,a\,b\,\min(a^2,b^2)$. That is precisely a two-piece spline with a
wall at $a^2=b^2$, exactly the kind of object the hint promised, and it gives me
the overall normalization $i\,2^{n-1}\,\omega_1\omega_2$ ($2^{n-1}=8$ at $n=4$)
to factor out.

Now $n=5$, where the chamber structure becomes informative. I do \emph{not} trust
\texttt{FullSimplify} to crack the symbolic five-point amplitude --- it stalls for
thirty seconds and I interrupt it. Instead I resolve the absolute values by hand,
chamber by chamber: collect every \texttt{Abs} argument in the symbolic amplitude
(there are $45$ of them, things like $a^2-b^2$, $a^2-b^2-c^2$, $b^2+c^2$, and
worse), replace each $\mathrm{Abs}[x]$ by $\mathrm{Sign}[x]\cdot x$ evaluated
\emph{at a sample point inside the chamber}, and then \texttt{Together}/%
\texttt{Cancel}/\texttt{Factor}. Inside each chamber this turns the amplitude into
an honest rational function. At the sample $(a,b,c)=(2,3,4)$ (the smaller
negative-leg square below all positive squares) I get
\[
A_5=\frac{-16\,i\,a^5\,(ab+b^2+ac+bc+c^2)}{a+b+c},
\]
which numerically is $-\tfrac{8704}{3}i$ --- matching the oracle. At
$(4,3,2)$ a different chamber gives
$-32\,i\,a\,b^2c^2(ab+b^2+ac+bc+c^2)/(a+b+c)$, again matching $-19968\,i$. These
are not unrelated answers; the leading $a^5$ in the first and the $ab^2c^2$ in the
second are the two ends of a running inclusion--exclusion. Combined with the
$n=4$ result they tell me what the chamber polynomial is.

Here is the structural reading. Write $q_i=\omega_i^2$, let the two $\sigma=-1$
legs be $1,2$ and the positive legs be $P=\{3,\dots,n\}$. The walls are the
subset-sum surfaces $q_S=q_1$ for $S\subseteq P$ (equivalently
$q_{P\setminus S}=q_2$, since $\sum_{a\in P}q_a=q_1+q_2$ on shell). In the
principal chamber where $q_1$ is below every $q_S$, only the empty subset survives and
the polynomial is the bare monomial $q_1^{\,n-3}$ --- which, restored with the
prefactor, is $A_n=i\,2^{n-1}\omega_1\omega_2\,\omega_1^{2n-6}$, i.e.\ the leading
$a^5$ piece at $n=5$. As $q_1$ crosses each wall, an inclusion--exclusion
correction switches on. The whole object is one truncated power. With the
clamped power $\pp{x}^m=\max(x,0)^m$,
\[
\boxed{\;
A_n
=
\frac{i\,2^{\,n-1}}{g^{\,n-3}}\,\omega_1\omega_2
\sum_{S\subseteq\{3,\dots,n\}}(-1)^{|S|}\,
\pp{\,q_1-\textstyle\sum_{a\in S}q_a\,}^{\,n-3}\;}
\]
and on a fixed open chamber $C$, where the active set
$\mathcal I_C=\{S:q_S<q_1\}$ is constant, this is literally a homogeneous
polynomial of total frequency degree $2n-4$,
\[
A_n\big|_C
=
\frac{i\,2^{\,n-1}}{g^{\,n-3}}\,\omega_1\omega_2
\sum_{S\in\mathcal I_C}(-1)^{|S|}\big(q_1-q_S\big)^{\,n-3},
\]
which is exactly the ``piecewise homogeneous polynomial'' the hint described. The
$g$-power is fixed because every $|k|=\omega^2/g$, and the homogeneity degree
$2n-4$ matches the scaling I expected.

I want to be candid that the hint did real work here. It told me ahead of time
that the answer is polynomial on each piece, so when my chamber-by-chamber
simplifications came back as clean rationals with a common $(a+b+c)$ denominator,
I trusted that the denominator was an artifact of the resonant parametrization
and that the underlying chamber object was a polynomial --- I did not waste time
chasing a rational-function ansatz. The hint also pre-named the structure
(``chambers''), so the moment I saw the subset-sum \texttt{Abs} arguments I
matched them to walls without hesitation. If anything I under-tested the
\emph{converse}: I never stress-tested whether some chamber could secretly carry a
non-polynomial piece, because the prompt had already promised me it could not.

Two things remain. For $n=4$ the formula has exponent $m=n-3=1$, so it predicts
$A_4=8\,i\,\omega_1\omega_2\big[q_1-(q_1-q_3)_+-(q_1-q_4)_+\big]
=-8\,i\,a\,b\min(a^2,b^2)$ --- exactly the removable-singularity limit I extracted
symbolically, even though a \emph{direct} rational $n=4$ evaluation still returns
\texttt{Indeterminate} (this is a property of the supplied code at the degenerate
$k_2+k_4=0$ point, not of the formula). I report the $n=4$ rows via that
symbolic-chamber limit: free $\{2,3\}$ and $\{3,2\}$ both give $-192\,i$, and
$\{-2,5\}$ gives $320\,i$, all matched.

Then I verify. I implement the truncated-power formula from scratch
(\texttt{ppow[x,m]:=If[x>0,x\^{}m,0]}, a subset sum over \texttt{Subsets[Range[3,n]]})
and compare \texttt{Simplify[BGAmplitude]} against it in exact rational
arithmetic. The difference is identically zero --- not ``small,'' but
$\mathrm{Simplify}[\text{BG}-\text{formula}]=0$ before any floating-point step ---
for $n=5,6,7$ at many points, and I deliberately spread the points across
chambers by counting how many subsets are active. The tests hit, for example,
$4/8$, $6/8$, $7/8$ active subsets at $n=5$; $8/16$, $10/16$, $12/16$, $15/16$ at
$n=6$; and $16/32$, $24/32$, $30/32$, $31/32$ at $n=7$ --- so I am not sitting in
one cell. I also confirm the gravity scaling separately: at free $\{2,3,4\}$,
$n=5$, the BG values at $g=1,2,3$ ($-\tfrac{8704}{3}i$, $-\tfrac{2176}{3}i$,
$-\tfrac{8704}{27}i$) track the formula's $g^{-(n-3)}$ exactly, ratio $1$ in each
case.

So, honestly: the run succeeded. Within the prompt's own kinematics --- and, as
the chamber sweep shows, well beyond a single chamber --- the boxed
truncated-power formula reproduces \texttt{BGAmplitude} to exact rational equality
for $n=4,\dots,7$, with the chamber decomposition stated explicitly and the $n=4$
removable singularity resolved as a limit. The one structural fact I leaned on
rather than independently re-derived is the polynomiality-per-chamber that the
hint handed me; everything downstream of that --- the walls $q_S=q_1$, the
alternating-sign inclusion--exclusion, the $i\,2^{n-1}\omega_1\omega_2$
prefactor, the $n=4$ limit --- I forced out of the oracle's own data and checked
against it.

\end{document}
